\subsection{Анализ затрат на изготовление продукта}

Затраты на выполнение проекта состоят из затрат на заработную плату исполнителям, затрат на закупку или аренду оборудования, затрат на организацию рабочих мест и затрат на накладные расходы.

Проектные затраты вычисляются согласно формуле:

\begin{equation}
K = C_{\text{ЗАРП}} + C_{\text{ОБ}} + C_{\text{ОРГ}} + C_{\text{НАКЛ}},
\end{equation}

где $C_{\text{ЗАРП}}$ -- заработная плата исполнителей;

$C_{\text{ОБ}}$ -- затраты на обеспечение необходимым оборудованием;

$C_{\text{ОРГ}}$ -- затраты на организацию рабочих мест;

$C_{\text{НАКЛ}}$ -- накладные расходы.

\subsubsection{Затраты на заработную плату}

Затраты на заработную плату исполнителям определяются формулой:

\begin{equation}
C_{\text{ЗАРП}} = C_{\text{ОСН}} + C_{\text{ДОП}} + C_{\text{ОТЧ}},
\end{equation}

где $C_{\text{ОСН}}$ -- основная заработная плата;

$C_{\text{ДОП}}$ -- дополнительная заработная плата;

$C_{\text{ОТЧ}}$ -- отчисления с заработной платы.

Расчет основной заработной платы при дневной оплате труда проводится на основе данных по окладам и графику занятости исполнителей. Её расчет выполняется согласно формуле:

\begin{equation}
C_{\text{ОСН}} = T_{\text{ЗАН}} \cdot O_{\text{ДН}},
\end{equation}

где $T_{\text{ЗАН}}$ -- число дней, отработанных исполнителем проекта;

$O_{\text{ДН}}$ -- дневной оклад исполнителя.

При условии восьмичасового рабочего дня дневной оклад рассчитывается по формуле:

\begin{equation}
O_{\text{ДН}} = \frac{8 \cdot O_{\text{МЕС}}}{F_{\text{М}}},
\end{equation}

где $O_{\text{МЕС}}$ -- месячный оклад,

$F_{\text{М}}$ -- месячный фонд рабочего времени.

Месячный оклад вычисляется с учетом налога на доходы физических лиц согласно формуле:

\begin{equation}
O_{\text{МЕС}} = O \cdot \left(1 + \frac{Н_{\text{НДФЛ}}}{100}\right),
\end{equation}

где $O$ -- "чистый" оклад,

$Н_{\text{НДФЛ}}$ -- налог на доходы физических лиц, составляющий 13\% от "чистый" оклада.

Месячный фонд рабочего времени принят на основании производственного календаря на 2026 год. При 40-часовой рабочей неделе годовая норма рабочего времени составляет 1972 часа, следовательно,

\begin{equation}
F_{\text{М}} = \frac{1972}{12} = 164.3 \ \text{час./мес.}
\end{equation}

После составления графика работ необходимо определить уровни заработной платы исполнителей $C_{\text{ОСН}}$. Для этого использованы данные HeadHunter. Согласно им, средний "чистый" месячный оклад инженера принят равным 125000 рублей, а аналитика данных -- 104400 рублей. Тогда месячный и дневной оклад для инженера составляют соответственно 141250 рублей и 6876.3 рубля, а для аналитика -- 117972 рубля и 5743.1 рубля.

Подробные результаты представлены в таблице~\ref{tab:salary_calc}.

\begin{table}[h!]
    \centering
    \caption{Расчет затрат на заработную плату исполнителей}
    \label{tab:salary_calc}
    \begin{tabular}{|p{3.2cm}|c|c|c|c|c|}
        \hline
        Исполнитель & $T_{\text{ЗАН}}$, дн. & $O$, руб. & $O_{\text{МЕС}}$, руб. & $O_{\text{ДН}}$, руб./дн. & $C_{\text{ОСН}}$, руб. \\
        \hline
        Инженер & 97.5 & 125000 & 141250 & 6876.3 & 670436.1 \\
        \hline
        Аналитик данных & 31.0 & 104400 & 117972 & 5743.1 & 178034.8 \\
        \hline
    \end{tabular}
\end{table}

Дополнительно произведём уточнение затрат на заработную плату на основе расчётов в КСП-системе.

Объём затрат на основную заработную плату в КСП-системе для инженера составил 173400{,}66 рубля за 228{,}8 часов работы, а для аналитика — 103189{,}63 рубля за 166{,}4 часа работы. Тогда суммарная основная заработная плата равна:
\begin{equation}
C_{\text{ОСН}} = 173400{,}66 + 103189{,}63 = 276590{,}29 \ \text{руб.}
\end{equation}

Дополнительные расходы на заработную плату возникают в результате выплат сотрудникам за время, проведённое вне рабочего процесса (например, отпуск, компенсации и т.д.). Методика расчёта дополнительной заработной платы определяется по формуле:
\begin{equation}
C_{\text{ДОП}} = 0.02 \cdot C_{\text{ОСН}}.
\end{equation}

В нашем случае:
\begin{equation}
C_{\text{ДОП}} = 0.02 \cdot 276590{,}29 = 5531{,}81 \ \text{руб.}
\end{equation}

Согласно Налоговому кодексу РФ, отчисления с заработной платы рассчитываются по формуле:
\begin{equation}
C_{\text{ОТЧ}} = N_{\text{СОЦ}} \cdot (C_{\text{ОСН}} + C_{\text{ДОП}}),
\end{equation}

где $N_{\text{СОЦ}}$ — страховые взносы, составляющие 30\% от выплат и включающие:
\begin{itemize}
    \item пенсионное страхование — 22\%;
    \item социальное страхование — 2{,}9\%;
    \item медицинское страхование — 5{,}1\%.
\end{itemize}

Тогда:
\begin{equation}
C_{\text{ОТЧ}} = 0.3 \cdot (276590{,}29 + 5531{,}81) = 84636{,}63 \ \text{руб.}
\end{equation}

Общие затраты на заработную плату исполнителям составляют:
\begin{equation}
C_{\text{ЗАРП}} = C_{\text{ОСН}} + C_{\text{ДОП}} + C_{\text{ОТЧ}} = 
276590{,}29 + 5531{,}81 + 84636{,}63 = 366758{,}73 \ \text{руб.}
\end{equation}